\documentclass[10pt,a4paper]{amsart}
\usepackage[margin=19mm]{geometry}
\usepackage{amsmath,amssymb,mathtools}
\usepackage[scr=rsfs,cal=euler]{mathalfa}
\usepackage{xfrac}
\usepackage[unicode,hidelinks,bookmarks=false]{hyperref}
\newcommand{\ie}{i.e.\ }
\newcommand{\cf}{cf.\ }
\newcommand{\resp}{respectively\ }
\newcommand{\Subsection}{\subsection}
\providecommand{\nobreakdash}{\nobreak\hbox{-}}
\setlength{\emergencystretch}{2em}
\multlinegap0pt
\usepackage{braket}
\usepackage{xcolor}
\usepackage{amssymb}
\newtheorem{lemma}{Lemma}
\newtheorem{thm}{Theorem}
\def\bbra#1{\langle\!\langle #1 \vert}
\def\bra#1{\langle #1 \vert}
\def\iinner#1#2{\langle\!\langle #1|#2\rangle\!\rangle} 
\def\ket#1{\vert #1 \rangle}
\def\kket#1{\vert #1 \rangle\!\rangle}
\title{Generating functions of dual $K$-theoretic $P$- and $Q$-functions and boson-fermion correspondence}
\author{Shinsuke Iwao}
\date{Source: arXiv:2301.12741v4 (CC BY 4.0)}
\begin{document}
\maketitle

\noindent\emph{Source and licence.} Generating functions of dual $K$-theoretic $P$- and $Q$-functions and boson-fermion correspondence. Version arXiv:2301.12741v4 (CC BY 4.0), distributed by arXiv under the Creative Commons Attribution 4.0 International licence, http://creativecommons.org/licenses/by/4.0/, as recorded in the arXiv licence metadata for this exact version and shown on the abstract page https://arxiv.org/abs/2301.12741v4. Full licence notice: \texttt{licenses/arxiv-2301.12741-CC-BY-4.0.txt}.\par\smallskip

\noindent\textit{Editorial scope.} Counted unit: the article's thm thm:fermion\_of\_gp (source0.tex lines 1572--1579). The article's complete original proof of it is delivered with it (source0.tex lines 1738--1830, one \texttt{proof} environment of 93 lines, which the article places away from the statement). No mathematics is written, repaired or re-derived: the statement and the proof are the article's own lines, in the article's own notation, and no retained line is substituted or re-flowed.  Retained uncounted as the source-local context the proof needs: lines 506--514; lines 1660--1672; lines 1674--1681.  Nothing was omitted from any retained span, so the package carries no omission record.  No retained line contains a citation, so the wrapper carries no bibliography and no bibliography entry.  No retained line contains a figure environment, an image-inclusion command or drawing code, so no image is delivered and the package declares no asset dependency.  Editorial formatting only, in addition: an AMS \texttt{amsart} preamble that restates the article's own declarations and macros used by the retained lines, namely the environments \texttt{lemma}, \texttt{thm} on separate counters, together with the standard packages those lines need.  The retained environments print in this document as Lemma 1, Theorem 1 in order of appearance, whereas the article prints the same items with the numbering of its own full document.  The complete native source of the article as distributed in the arXiv e-print for this exact version is bundled unchanged as \texttt{sources/136679/source0.tex}.
\begin{lemma}[Duality relation]
\label{lemma:duality}
We have the following commutation relations:
\[
[(\Phi^{(\beta)}_m)^\ast, \phi^{[\beta]}_n]_+
=[(\Phi^{[\beta]}_m)^\ast, \phi^{(\beta)}_n]_+
=\delta_{m,n}.
\]
\end{lemma}

\begin{equation}\label{eq:mokuhyou}
%\iinner{\mu}{\lambda}=
%\begin{cases}
%1 & (\lambda=\mu\text{ and }r \text{ is even})\\
%1 & ((\lambda,0)=\mu\text{ and }r \text{ is odd}),\\
%0 & (\text{otherwise}).
%\end{cases}
\iinner{\mu}{\lambda}=
\begin{cases}
1 & (\mu=\lambda \text{ or } \mu=(\lambda,0)),\\
0 & (\text{otherwise}).
\end{cases}
\end{equation}

\begin{lemma}\label{lemma:hojyo2}
We have the following relations:
\begin{enumerate}
\item[(A).] When $\mu=\emptyset$ and $n\geq 0$, we have $\bbra{\emptyset}(\Phi^{[\beta]}_n-\tfrac{1}{2}(-\tfrac{\beta}{2})^{n})=0$.
\item[(B).] When $\mu\neq \emptyset$ and $n> \mu_1$, we have $\bbra{\mu}(\Phi^{[\beta]}_n-\tfrac{1}{2}(-\tfrac{\beta}{2})^{n})=0$.
\item[(C).] When $\lambda\neq \emptyset$ and $m>\lambda_1$, we have $(\phi^{(\beta)}_m)^\ast \kket{\lambda}=0$.
\end{enumerate}
\end{lemma}

\begin{thm}\label{thm:fermion_of_gp}
For a strict partition $\lambda=(\lambda_1>\dots>\lambda_r>0)$, we have
\begin{equation}\label{eq:def_of_gp}
\begin{aligned}
gp_\lambda(x)=\bra{0}e^{\mathcal{H}^{[\beta]}}\ket{\lambda}_p.
\end{aligned}
\end{equation}
\end{thm}

\begin{proof}[Proof of Theorem \ref{thm:fermion_of_gp}]

Let $E=\iinner{\mu}{\lambda}$.
To prove the theorem, it suffices to verify \eqref{eq:mokuhyou}.
We show \eqref{eq:mokuhyou} by induction on $s\geq 0$.
We begin by considering the following base cases:
\begin{equation}\label{eq:base_case}
\iinner{\mu}{\emptyset}=
\begin{cases}
1 & (\mu=\emptyset),\\
1 & (\mu=(0)),\\
0 & (\text{otherwise}),\\
\end{cases}\qquad
\iinner{\emptyset}{\lambda}
=
\begin{cases}
1 & (\lambda=\emptyset),\\
0 & (\lambda\neq \emptyset).
\end{cases}
\end{equation}

The first equation in \eqref{eq:base_case} can be verified as follows: When $\mu=\emptyset$, we have $E=\bra{0}(\phi_0+1)\ket{0}=1$.
When $\mu=(0)$, we have \[
E=
\bra{0}(\phi_0+1)e^\theta(\phi_0^{(\beta)})^\ast\ket{0}
=\bra{0}(\phi_0+1)e^\theta\phi_0\ket{0}
=\bra{0}(\phi_0^2+\phi_0)\ket{0}=1.
\]
When $\mu\neq \emptyset$ and $\mu\neq(0)$, we have $E=0$ from Lemma \ref{lemma:hojyo2} (C).

The second equation in \eqref{eq:base_case} can be verified as follows: We have already shown $E=1$ when $\lambda=\emptyset$.
When $\lambda\neq \emptyset$, we have $E=0$ from Lemma \ref{lemma:hojyo2} (A).

We proceed for the case when $s,r>0$.
% we prove \eqref{eq:mokuhyou} by induction on $s$.
If $\mu_1<\lambda_1$, then $E=0$ by Lemma \ref{lemma:hojyo2} (B).
If $\mu_1>\lambda_1$, then $E=0$ by Lemma \ref{lemma:hojyo2} (C).
If $\mu_1=\lambda_1(>0)$, then we have
\[
\begin{aligned}
E
&=
\bbra{\mu'}e^{\theta}(\phi^{(\beta)}_{\mu_1})^\ast
(\Phi^{[\beta]}_{\lambda_1}-\tfrac{1}{2}(-\tfrac{\beta}{2})^{\lambda_1})e^{-\theta}\kket{\lambda'}\\
&=
\bbra{\mu'}e^{\theta}
\left[
(\phi^{(\beta)}_{\mu_1})^\ast,
\Phi^{[\beta]}_{\lambda_1}
\right]_+
e^{-\theta}\kket{\lambda'}
-
\bbra{\mu'}e^{\theta}
\Phi^{[\beta]}_{\lambda_1}
(\phi^{(\beta)}_{\mu_1})^\ast
e^{-\theta}\kket{\lambda'}
%\\
%&\hspace{15em}
-\tfrac{1}{2}(-\tfrac{\beta}{2})^{\lambda_1}
\bbra{\mu'}e^{\theta}
(\phi^{(\beta)}_{\mu_1})^\ast
e^{-\theta}\kket{\lambda'}\\
&
=
\iinner{\mu'}{\lambda'}
-
\bbra{\mu'}e^{\theta}
\Phi^{[\beta]}_{\lambda_1}
(\phi^{(\beta)}_{\mu_1})^\ast
e^{-\theta}\kket{\lambda'}
-\tfrac{1}{2}(-\tfrac{\beta}{2})^{\lambda_1}
\bbra{\mu'}e^{\theta}
(\phi^{(\beta)}_{\mu_1})^\ast
e^{-\theta}\kket{\lambda'}
\qquad 
(\mathrm{Lemma\,\ref{lemma:duality}}).
\end{aligned}
\]
The last two terms in the last expression equal to $0$ because
\[
(\phi^{(\beta)}_{\mu_1})^\ast
e^{-\theta}\kket{\lambda'}=
e^{-\theta}
(\phi^{(\beta)}_{\mu_1}-\beta \phi^{(\beta)}_{\mu_1+1}+\cdots)^\ast
\kket{\lambda'}
=
0
\]
by Lemma~\ref{lemma:hojyo2}\,(C).
Hence, we obtain $E=\iinner{\mu'}{\lambda'}$, where
$\mu'=(\mu_2>\dots>\mu_s\geq 0)$ and $\lambda'=(\lambda_2>\dots>\lambda_r>0)$.
By induction hypothesis, we conclude \eqref{eq:mokuhyou}.
\end{proof}

\end{document}
